\section{고찰}

\subsection{선행연구와의 비교}
\begin{enumerate}
\item \textbf{원리.} 결과 체인링은 모멘트팔이 \(\theta\approx130^\circ\), \(310^\circ\)에서 가장 크고, 식 (\ref{eq:1})에 따라 이때 크랭크가 평균보다 23\% 느리게 돈다. \(310^\circ\)는 오른다리가 앞쪽에서 내려가며 미는 구간(\(221.5^\circ\)에서 \(399^\circ\))의 가운데이므로, 힘을 내는 행정에서 크랭크를 늦춘다는 Rankin과 Neptune의 결과\cite{rankin2008}와 같다.
\item \textbf{크기.} 그들의 최적 지름비는 1.24--1.35로 조건 안쪽에서 정해졌지만, 본 연구의 최적은 조건으로 둔 상한 2에 닿았다. 그들의 모형에는 근육이 켜지고 꺼지는 데 걸리는 시간이 있어 지나치게 찌그러진 모양이 손해를 보지만, 본 모델에는 이것이 없기 때문으로 보인다. 효과의 크기(약 3\%)는 비슷하다.
\item \textbf{효과 지수.} 케이던스가 높을수록 IE가 작아지는 경향은 선행연구\cite{bini2013}와 같다. 반면 선행연구에서는 출력이 높을수록 IE가 커지는데, 본 모델에서는 모든 힘이 출력에 비례해 IE가 출력과 상관없다. 다리 무게와 근육의 반응 시간이 모델에 없기 때문이다.
\item \textbf{관절 분담.} 고관절이 주로 일하는 것은 60 rpm 이하이다\cite{aasvold2019}. 고관절만 힘을 내는 본 모델의 결과 중 60 rpm 값(1.63\%)이 현실에 더 가깝고, 무릎이 주로 일하는 100 rpm 값(2.99\%)은 실제 효과를 크게 볼 수 있다.
\item \textbf{실제 사람.} 측정 연구에서는 비원형 체인링이 최대 출력을 바꾸지 못했고, 사람은 발목 움직임으로 무릎과 고관절 움직임을 유지하였다\cite{leong2017}. 본 연구의 결과는 최대 출력이 아니라 같은 출력에서의 고관절 부담에 관한 것이므로 직접 비교는 어렵다.
\end{enumerate}

\subsection{모양을 한 가지 꼴로 정하지 않은 의미}
이론 문서의 \(\cos2\theta\) 한 성분 꼴은 볼록 조건만 두면 \(e=1/3\), 지름비 2에서 효과가 가장 크지만, 그때 긴 축 끝이 뾰족해 실제로 쓸 수 없다. 모양을 푸리에 급수로 넓히고 볼록 최적화로 정확히 풀면, 제작 조건 안에서 무엇이 가능한지가 꼴의 선택과 상관없이 정해진다. 그 결과 \(\cos4\theta\) 성분 하나가 꼭 필요하고, 그 이상은 이득이 작다는 것을 알 수 있었다. 타원은 이 조건을 자연스럽게 만족하는 간단한 모양이다.

\subsection{한계}
\begin{itemize}
\item 이상적 배분에서는 올라가는 다리도 약 150 N으로 페달을 당긴다([Fig.~\ref{fig:clock}]). 선수에게 IE를 높이라고 지시하면 효율이 떨어진다는 결과\cite{bini2013}를 고려하면, 이상적 배분은 실제보다 당기는 힘을 크게 잡은 경우이다.
\item 무릎과 발목이 힘을 내지 않는다는 가정, 장력 조절 장치의 \(T_-\)를 0으로 둔 가정이 결과의 크기에 영향을 준다.
\item 지름비 상한(2)과 최소 곡률반지름(25 mm)은 실제 변속 장치와 체인 조건에서 정한 값이 아니다. 결과는 지름비 상한에 닿아 있으므로 이 값에 따라 효과가 달라진다.
\end{itemize}

\section{결론}
비원형 체인링의 운동방정식을 정의와 법칙에서 유도하고, 고관절만 힘을 내는 다리 모델로 한 바퀴 중 가장 큰 고관절 부담을 최소화하는 체인링 모양을 찾았다. 모양을 푸리에 급수로 일반화하여 볼록 최적화로 정확히 풀었고, 다음을 얻었다.
\begin{enumerate}
\item 최적 모양은 \(180^\circ\) 대칭이며, 제약이 없으면 \(\rho_+\propto\sqrt K\)이다. 이때 부담 감소의 상한은 60 rpm에서 2.60\%, 100 rpm에서 4.72\%이다.
\item 지름비 2, 곡률반지름 25 mm 조건에서 타원형 체인링은 가장 큰 부담을 60 rpm에서 1.63\%, 100 rpm에서 2.99\% 줄이고, 위쪽 체인 장력은 한 바퀴 내내 양수이다. 푸리에 \(M=24\)까지 넓혀도 추가 이득은 0.15\%p 이하이다.
\item 무릎이 힘을 내지 않으면 페달 힘은 정강이 방향으로만 전달되며, 이 힘의 분해로 구한 효과 지수는 원형 0.583에서 0.654로 커진다(60 rpm). 그 원인은 데드스팟 근처에 머무는 시간이 줄어든 데 있다.
\end{enumerate}

향후에는 무릎에 토크를 하나 더 넣은 다리 모델, 근육이 켜지고 꺼지는 시간을 넣은 근육 모델, 실측 페달 힘 곡선을 배분으로 쓰는 계산으로 결과의 크기를 다시 확인하고, 실제 체인과 변속 조건에서 지름비와 곡률반지름의 범위를 정할 계획이다.
